\chapter{Installation and requirements}\label{ch:install}
\chapterintro{AgenticTO is a Python package with a web front end and a separate GPU solver. This chapter lists what you
need, installs the Python environment and the front end, connects the \code{voxfem} solver, and sets up the optional
language agents, which run on personal subscriptions only.}

\section{Requirements}\label{sec:requirements}

\begin{table}[h]
\caption{Requirements. ``Tested with'' is the configuration used for every result in the paper.}\label{tab:req}
\begin{fixtable}{@{}>{\raggedright\arraybackslash}p{0.2\linewidth}>{\raggedright\arraybackslash}p{0.27\linewidth}X@{}}
Item & Tested with & Notes \\ \midrule
GPU & NVIDIA GeForce RTX 4090, 24 GB & A CUDA 12 GPU. GPU memory bounds the finest grid you can release; rule
\rref{R013} estimates the memory before a solve starts. \\
Host & Intel Core i9-13900KF, 96 GB RAM & Large releases keep density fields in host memory as well. \\
Operating system & Windows 11 & The spill guard reads the Windows GPU-memory counters. Other platforms have not been
tested. \\
Python & 3.10 or newer (a conda environment) & See Section~\ref{sec:pyenv} for packages. \\
Node.js & 22 & Only for the web front end. \\
Solver & \code{voxfem} & Included in \code{src/voxfem/} (BSD 3-Clause; Section~\ref{sec:voxfem}). \\
Browser & Chrome & The workspace has been tested in Chrome with one user per session. \\
Codex CLI (optional) & 0.165.0, ChatGPT login & Runs the built-in assistant. \\
Claude Code (optional) & 2.1.312, claude.ai login & An external agent over MCP. \\
\end{fixtable}
\end{table}

\begin{notebox}[Shared desktop GPU]
In the paper the GPU also drove the desktop, and other processes held about 2--5~GB of its memory. That is why the
default ladder stops below \LadderMax\ elements, although a 24~GB card could hold more on its own. A guarded attempt at
$1.13\times10^{8}$ elements was stopped twice by the spill guard and is not part of any reported result
\claim{C3.2}.
\end{notebox}

\section{The Python environment}\label{sec:pyenv}
Create a conda environment with a CUDA 12 build of CuPy and the packages AgenticTO imports:

\begin{shell}
conda create -n agentto python=3.11
conda activate agentto
pip install numpy scipy matplotlib "cupy-cuda12x" "pydantic>=2.8,<3" \
    fastapi uvicorn mcp vtk pillow websockets pytest
\end{shell}

What each package is for: \code{numpy}, \code{scipy} and \code{cupy} for the numerics and \code{field\_stats};
\code{pydantic} for the specification schema and the tool arguments; \code{fastapi}, \code{uvicorn} and
\code{websockets} for the backend; \code{mcp} for the agent interface; \code{matplotlib} for the setup pictures of the
\tool{render} tool; \code{vtk} for the marching-cubes surface and the STL export; \code{pillow} for thumbnails.

Optional extras:
\begin{itemize}
  \item \code{pywin32} on Windows. The spill guard then reads the GPU memory counters in-process (about 0.1~ms per
  sample) instead of through a PowerShell subprocess (about 2.4~s per sample).
  \item \code{pyvista} for the 3D renders of the paper figures, \code{scikit-fem} for the independent cross-check of
  study X0, and \code{openpyxl} for the adjudication sheets (Chapter~\ref{ch:reproduce}).
\end{itemize}

AgenticTO is not installed into the environment. The scripts put \code{src/} on the Python path themselves, and the
MCP configurations in Chapter~\ref{ch:mcp} set \code{PYTHONPATH}. If you prefer an installed package,
\code{pip install -e .} in the repository root uses \code{pyproject.toml}.

\section{The web front end}
\begin{shell}
cd web
npm install
cd ..
\end{shell}
This installs React, three.js and Vite at the versions pinned in \code{web/package.json}.

\section{Connecting the solver}\label{sec:voxfem}
\index{voxfem@\code{voxfem}}\index{AGENTTO_VOXFEM_SRC@\code{AGENTTO\_VOXFEM\_SRC}}
AgenticTO runs on \code{voxfem}, the GPU solver of the authors' companion work on fused matrix-free operators,
Galerkin multigrid and guarded state solves. \code{voxfem} ships with the AgenticTO repository in \code{src/voxfem/}, unchanged and
under its own BSD 3-Clause License, and is found automatically when \code{src/} is on the path. AgenticTO imports it read-only, never writes into
it (no bytecode caches either), and extends it only by subclassing.

To use another copy of the solver, point AgenticTO at the directory that contains that \code{voxfem} package (the
directory you would put on \code{PYTHONPATH} to make \code{import voxfem} work):

\begin{shell}
# POSIX shell
export VOXFEM_SRC=/path/to/other/src
# Windows PowerShell
$env:VOXFEM_SRC = "D:/other/src"
\end{shell}

Set the variable in the shell that starts the backend, the MCP server or a script; the GPU worker inherits it.
The older name \code{AGENTTO\_VOXFEM\_SRC} still works.

\begin{warnbox}[Without a GPU]
Without a CUDA GPU you can still read the code, the briefs, the logs and every result file, and regenerate the
paper's tables and data figures from those files (Chapter~\ref{ch:reproduce}). Every solve needs a CUDA GPU.
\end{warnbox}

\section{Checking the installation}
\begin{shell}
python -B -m pytest -q
\end{shell}
The tests that solve on the GPU need a free GPU. Run one GPU job at a time: another CUDA program on the
same card can push a solve into shared memory, and the spill guard will then stop it.

\section{Optional: language agents on a subscription}\label{sec:agents}
\index{subscription policy}\index{PolicyError}
AgenticTO calls language models only through command-line tools that are logged in with a personal subscription:
the Codex CLI with a ChatGPT login for the built-in assistant, and optionally Claude Code with a claude.ai login.

\begin{shell}
codex login        # sign in with your ChatGPT account
claude             # the first run asks you to sign in with your claude.ai account
\end{shell}

Before every call AgenticTO checks two things, and refuses to start the agent if either fails:
\begin{enumerate}\sloppy
  \item No API-key variable is set: \code{OPENAI\_API\_KEY}, \code{CODEX\_API\_KEY}, \code{ANTHROPIC\_API\_KEY},
  \code{ANTHROPIC\_AUTH\_TOKEN}, \code{CLAUDE\_CODE\_USE\_BEDROCK} or \code{CLAUDE\_CODE\_USE\_VERTEX}. Any of them would
  silently switch billing to metered API use.
  \item The CLI is logged in with a subscription: \code{codex login status} must report ChatGPT, and
  \code{claude auth status} must report \code{claude.ai}.
\end{enumerate}
A refusal appears in the assistant panel as a \code{PolicyError}. Unset the variable, or log in again, and retry. The
child processes are started with those variables removed from their environment as a second safeguard.

The built-in assistant uses the model \code{gpt-6.1-sol} at reasoning effort high (the defaults in
\code{src/agentto/llm/cli.py}). Set \code{AGENTTO\_CODEX} or \code{AGENTTO\_CLAUDE} if the CLI is not on your
\code{PATH}.

\begin{tipbox}[The assistant is optional]
Every design in the paper's 50-brief suite can be built and solved in the workspace alone. The assistant and external
agents use the same tools you do.
\end{tipbox}
